\section{The Order of an Element and Primitive Roots}

\begin{definition}[Order of an Element]
Let $p$ be a prime and $a \in \mathbb{Z}$ so that $p$ does not divide $a$. The order of $a$ modulo $p$, denoted $\text{ord}_p(a)$, is the minimum positive integer $n \in \mathbb{N} \setminus \{0\}$ so that $a^n \equiv 1 \pmod{p}$.
\end{definition}

\begin{center}
    \vspace{10pt}
    \line(1,0){450}
    \vspace{25pt}
\end{center}

\begin{proposition}[Order Divisibility Property]
Let $p$ be a prime, $a \in \mathbb{Z}$ so that $p$ does not divide $a$, and $n \ge 1$. Then $a^n \equiv 1 \pmod{p}$ if and only if $\text{ord}_p(a)$ divides $n$.
\end{proposition}
\begin{proof}
We demonstrate both directions of the biconditional.

\textbf{First Direction ($\impliedby$)}: Assume $\text{ord}_p(a) \mid n$.
$\rightarrow \exists k \in \mathbb{Z}$ so that $n = k \cdot \text{ord}_p(a)$.
Then:
$$
a^n = a^{k \cdot \text{ord}_p(a)} = (a^{\text{ord}_p(a)})^k
$$
By definition of the order, $a^{\text{ord}_p(a)} \equiv 1 \pmod{p}$. Therefore:
$$
a^n \equiv (1)^k \equiv 1 \pmod{p}
$$

\textbf{Second Direction ($\implies$)}: Assume $a^n \equiv 1 \pmod{p}$.
By the division algorithm, $\exists q, r \in \mathbb{Z}$ so that $n = q \cdot \text{ord}_p(a) + r$, where $0 \le r < \text{ord}_p(a)$.
We evaluate $a^n$:
$$
a^n = a^{q \cdot \text{ord}_p(a) + r} = (a^{\text{ord}_p(a)})^q \cdot a^r
$$
Substituting the assumed and known congruences:
$$
1 \equiv (1)^q \cdot a^r \pmod{p} \rightarrow a^r \equiv 1 \pmod{p}
$$
Since $\text{ord}_p(a)$ is the *minimum* positive integer for which this holds, and $0 \le r < \text{ord}_p(a)$, the only possibility is that $r = 0$.
$\rightarrow \text{ord}_p(a)$ divides $n$.
\end{proof}

\begin{center}
    \vspace{10pt}
    \line(1,0){450}
    \vspace{25pt}
\end{center}

\begin{proposition}[Equality of Powers]
Let $p$ be a prime and $g \in (\mathbb{Z}/p\mathbb{Z})^*$. Then $g^x = g^y$ if and only if $x \equiv y \pmod{\text{ord}_p(g)}$.
\end{proposition}
\begin{proof}
Since $g \in (\mathbb{Z}/p\mathbb{Z})^*$, $g$ is invertible.
$$
g^x = g^y \iff g^x \cdot g^{-y} = 1 \iff g^{x-y} = 1
$$
Applying the proposition from the previous section (2.1), $g^n = 1$ if and only if $\text{ord}_p(g) \mid n$.
Let $n = x-y$. We conclude:
$$
g^{x-y} = 1 \iff \text{ord}_p(g) \mid (x-y)
$$
By the definition of modular congruence, this is equivalent to: $x \equiv y \pmod{\text{ord}_p(g)}$.
\end{proof}

\begin{center}
    \vspace{10pt}
    \line(1,0){450}
    \vspace{25pt}
\end{center}

\begin{definition}[Primitive Root]
An element $g \in (\mathbb{Z}/p\mathbb{Z})^*$ is a primitive root if its powers generate all the elements of $(\mathbb{Z}/p\mathbb{Z})^*$. That is, the set $\{ g^1, g^2, \ldots, g^{p-1} \}$ is equal to the set $\{ 1, 2, \ldots, p-1 \}$.
\end{definition}

\begin{center}
    \vspace{10pt}
    \line(1,0){450}
    \vspace{25pt}
\end{center}

\begin{theorem}[Primitive Root Condition]
An element $g \in (\mathbb{Z}/p\mathbb{Z})^*$ is a primitive root if and only if $\text{ord}_p(g) = p-1$.
\end{theorem}
\begin{proof}[Proof Outline]
For $g$ to be a primitive root, the $p-1$ powers $\{ g^1, g^2, \ldots, g^{p-1} \}$ must all be distinct.
By the Proposition on the Condition for Equality of Powers, $g^i = g^j$ if and only if $i \equiv j \pmod{\text{ord}_p(g)}$.
For the powers to be distinct for all $1 \le i < j \le p-1$, we must ensure that $g^{j-i} \ne 1$ for any exponent $j-i$ where $1 \le j-i < p-1$.
This means that the smallest positive exponent $k$ so that $g^k = 1$ must be $p-1$.
By definition, $\text{ord}_p(g) = k$.
Therefore, $g$ is a primitive root if and only if $\text{ord}_p(g) = p-1$.
\end{proof}

\bigskip
\begin{definition}
    Let $p$ be a prime, and $a \in \mathbb{Z}$ such that $p \nmid a$.
    The \textbf{order} $ord_{p}(a)$ of $a$ modulo $p$ is the minimum $n \in \mathbb{N}$ such that $n \geq 1$ and $a^{n} \equiv 1 \pmod{p}$.
\end{definition}

\bigskip
\begin{proposition}
    Let $p$ be a prime, $a \in \mathbb{Z}$ such that $p \nmid a$, and $n \ge 1$.
    Then $a^{n} \equiv 1 \pmod{p}$ if and only if the order of $a$ modulo $p$, denoted $\text{ord}_{p}(a)$, divides $n$.
    In particular, $\text{ord}_{p}(a) \mid (p-1)$ (a consequence of Fermat's Little Theorem).
\end{proposition}

\begin{proof}
    Let $k = \text{ord}_{p}(a)$.
    
    \textbf{$(\Rightarrow)$ If $a^{n} \equiv 1 \pmod{p}$, then $k \mid n$} \\
    Use the Division Algorithm on $n$ by $k$: $n = qk + r$, where $0 \le r < k$.
    
    $$1 \equiv a^{n} \equiv a^{qk + r} \equiv (a^{k})^{q} \cdot a^{r} \equiv 1^{q} \cdot a^{r} \equiv a^{r} \pmod{p}$$
    
    Since $a^{r} \equiv 1 \pmod{p}$ and $k$ is the \textbf{minimum positive exponent} satisfying this property, and $r < k$, we must have $r=0$.
    Thus, $n = qk$, which means $k \mid n$.

    \textbf{$(\Leftarrow)$ If $k \mid n$, then $a^{n} \equiv 1 \pmod{p}$} \\
    If $k \mid n$, then $n = qk$ for some $q \in \mathbb{Z}$.
    $$a^{n} = a^{qk} = (a^{k})^{q} \equiv 1^{q} \equiv 1 \pmod{p}$$
\end{proof}

\begin{center}
    \vspace{10pt}
    \line(1,0){450}
    \vspace{25pt}
\end{center}

\begin{corollary}
    Let $p$ be a prime, $a \in \mathbb{Z}$ such that $p \nmid a.$
    $ord_{p}(a) = \min\{n \in \mathbb{N} : n \geq 1 \text{ so that } a^{n} \equiv 1 \pmod{p}\}$
    $= \min\{n \in \mathbb{N} : n \geq 1 \text{ so that } a^{n} \equiv 1 \pmod{p} \text{ and } n \mid p-1\}$.
\end{corollary}

\begin{center}
    \vspace{10pt}
    \line(1,0){450}
    \vspace{25pt}
\end{center}

\begin{proposition}
    Let $p$ be a prime. If $g \in \mathbb{F}_{p}^{*}$, then $g^{x} = g^{y} \iff y \equiv x \pmod{ord_{p}(g)}$.
\end{proposition}
\begin{proof}
    ($\Leftarrow$) Assume $x \equiv y \pmod{ord_{p}(g)}$. Then $x-y = k \cdot ord_{p}(g)$ for some integer $k$.
    $g^{x} = g^{y+k \cdot ord_{p}(g)} = g^{y} \cdot (g^{ord_{p}(g)})^{k} \equiv g^{y} \cdot 1^{k} \equiv g^{y} \pmod{p}$.

    ($\Rightarrow$) Assume $g^{x} = g^{y}$. Without loss of generality, assume $x \geq y$.
    Since $g \in \mathbb{F}_{p}^{*}$ is invertible, multiplying by $(g^{-1})^{y}$ gives $g^{x-y} = 1$.
    By the previous Proposition, $ord_{p}(g) \mid x-y$, which means $x \equiv y \pmod{ord_{p}(g)}$.
\end{proof}

\begin{center}
    \vspace{10pt}
    \line(1,0){450}
    \vspace{25pt}
\end{center}

\begin{definition}
    Let $p$ be a prime.
    A \textbf{primitive root} of $\mathbb{F}_{p}$ is any $g \in \mathbb{F}_{p}^{*}$ such that the powers of $g$ generate all the elements of $\mathbb{F}_{p}^{*}$, i.e., $\mathbb{F}_{p}^{*} = \{1, g, g^{2}, \dots, g^{p-2}\}$.
\end{definition}

\begin{center}
    \vspace{10pt}
    \line(1,0){450}
    \vspace{25pt}
\end{center}

\begin{theorem}[Characterization of the Primitive Roots]
    Let $p$ be a prime. $g \in \mathbb{F}_{p}^{*}$ is a primitive root if and only if $ord_{p}(g) = p-1$.
\end{theorem}
\begin{proof}
    The proof establishes the equivalence between $g$ generating the group $\mathbb{F}_{p}^{*}$ and $g$ having the maximum possible order. Note that $|\mathbb{F}_{p}^{*}| = p-1$.

    \textbf{$(\Rightarrow)$ If $g$ is a primitive root, then $\text{ord}_{p}(g) = p-1$} \\
    By definition, if $g$ is a primitive root, the set of its powers $\{1, g, \dots, g^{p-2}\}$ contains \textbf{all $p-1$ distinct elements} of $\mathbb{F}_{p}^{*}$.
    
    Let $m = \text{ord}_{p}(g)$. If we assume $m < p-1$, then the sequence of powers $1, g, g^2, \dots$ repeats with a period of $m$. This means the set of distinct powers $\langle g \rangle$ has size $|\langle g \rangle| = m$.
    
    Since $g$ is a primitive root, $\langle g \rangle = \mathbb{F}_{p}^{*}$, so $m = |\mathbb{F}_{p}^{*}| = p-1$.
    The assumption $m < p-1$ leads to the contradiction $|\langle g \rangle| < |\mathbb{F}_{p}^{*}|$. Thus, $m = p-1$.

    \textbf{$(\Leftarrow)$ If $\text{ord}_{p}(g) = p-1$, then $g$ is a primitive root} \\
    Assume $\text{ord}_{p}(g) = p-1$. We must show the powers $1, g, \dots, g^{p-2}$ are all distinct.
    
    Suppose, for contradiction, that $g^{i} \equiv g^{j} \pmod{p}$ for $0 \le i < j \le p-2$.
    Since $g$ is invertible, $g^{j-i} \equiv 1 \pmod{p}$.
    
    Let $k = j-i$. Then $0 < k < p-1$.
    This implies that a power $g^k \equiv 1 \pmod{p}$ for an exponent $k$ smaller than $p-1$. This contradicts the assumption that $\text{ord}_{p}(g) = p-1$ is the \textbf{minimum} positive exponent that results in $1$.
    
    Thus, the $p-1$ powers $\{1, g, \dots, g^{p-2}\}$ are distinct. Since $\mathbb{F}_{p}^{*}$ has exactly $p-1$ elements, the powers must generate the entire group, meaning $g$ is a primitive root.
\end{proof}

\begin{center}
    \line(1,0){450}
\end{center}

\begin{theorem}[Existence of Primitive Roots]
    Let $p$ be a prime. $\mathbb{F}_{p}$ has at least one primitive root. (Admitted)
\end{theorem}

\bigskip

\begin{proof}
    Let $x = g^{\frac{p-1}{2}}$. We first observe the square of $x$:
    $$x^{2} = \left(g^{\frac{p-1}{2}}\right)^{2} = g^{p-1}$$
    
    By \textbf{Fermat's Little Theorem}, $g^{p-1} \equiv 1 \pmod{p}$, provided $p \nmid g$ (which is true since $g \in \mathbb{F}_p^*$).
    Thus, $x$ satisfies the congruence:
    $$x^{2} \equiv 1 \pmod{p}$$
    
    In the field $\mathbb{F}_{p}$, the only solutions to $x^{2} \equiv 1$ are $x=1$ and $x=-1$ (or $x \equiv p-1 \pmod{p}$).
    
    We must show that $x \neq 1$. We proceed by contradiction.
    
    Assume $x=1$, meaning:
    $$g^{\frac{p-1}{2}} \equiv 1 \pmod{p}$$
    
    By the property of the order of an element, this would imply:
    $$\text{ord}_{p}(g) \mid \frac{p-1}{2}$$
    
    However, since $g$ is a \textbf{primitive root}, we know by definition that $\text{ord}_{p}(g) = p-1$.
    The divisibility condition becomes:
    $$(p-1) \mid \frac{p-1}{2}$$
    
    Since $p>2$, this condition is false (a positive integer cannot divide a number half its size). This is a \textbf{contradiction}.
    
    Therefore, the only remaining possibility for the value of $x$ is $x=-1$.
    $$g^{\frac{p-1}{2}} \equiv -1 \pmod{p}$$
\end{proof}

\begin{center}
    \vspace{10pt}
    \line(1,0){450}
    \vspace{25pt}
\end{center}

\begin{remark}
    Let $p>2$ be a prime. If $g$ is a square in $\mathbb{F}_{p}$, then $g$ is not a primitive root in $\mathbb{F}_{p}$.
\end{remark}
\begin{proof}
    Let $g$ be a square, so $g = h^{2}$ for some $h \in \mathbb{F}_{p}^{*}$.
    Then $g^{(p-1)/2} = (h^{2})^{(p-1)/2} = h^{p-1} \equiv 1 \pmod{p}$ by Fermat's Little Theorem.
    Since $g^{(p-1)/2} \equiv 1 \pmod{p}$, by the proposition on order, $ord_{p}(g) \mid \frac{p-1}{2}$.
    Since $\frac{p-1}{2} < p-1$ (as $p>2$), we have $ord_{p}(g) < p-1$.
    By the characterization theorem, $g$ is not a primitive root.
\end{proof}

\begin{center}
    \vspace{10pt}
    \line(1,0){450}
    \vspace{25pt}
\end{center}

\begin{lemma}[Effective Test for a Primitive Root Candidate]
    Let $p>2$ be prime and let $p-1 = \prod_{i=1}^{m}p_{i}^{a_{i}}$ be the prime factorization of $p-1$.
    Then $g \in \mathbb{F}_{p}^{*}$ is a primitive root if and only if $g^{(p-1)/p_{i}} \not\equiv 1 \pmod{p}$ for all $i=1, \dots, m$.
\end{lemma}

\begin{proof}
    Let $k = \text{ord}_{p}(g)$. Recall that $g$ is a primitive root $\iff k = p-1$.
    
    \textbf{$(\Rightarrow)$ If $g$ is a primitive root, the test passes} \\
    Assume $g$ is a primitive root, so $k = p-1$.
    
    Suppose, for contradiction, that $g^{(p-1)/p_{i}} \equiv 1 \pmod{p}$ for some prime factor $p_i$.
    By the property of the order, $k$ must divide the exponent:
    $$k \mid \frac{p-1}{p_{i}} \implies (p-1) \mid \frac{p-1}{p_{i}}$$
    Since $p_i$ is a prime ($p_i \ge 2$), this divisibility is impossible.
    Thus, $g^{(p-1)/p_{i}} \not\equiv 1 \pmod{p}$ must hold for all $i$.

    \textbf{$(\Leftarrow)$ If the test passes, $g$ is a primitive root} \\
    Assume $g^{(p-1)/p_{i}} \not\equiv 1 \pmod{p}$ for all $i$.
    We know $k \mid p-1$, so $k$ is a divisor of $p-1$, and thus $k$ must be of the form $k = p_{1}^{f_{1}} \cdots p_{m}^{f_{m}}$ with $0 \le f_{i} \le a_{i}$.
    
    Suppose, for contradiction, that $k < p-1$. This means at least one exponent $f_i$ must be strictly less than $a_i$ (i.e., $f_i \le a_i - 1$).
    
    For such an index $i$, we can construct the exponent $\frac{p-1}{p_{i}}$. Since $p_i^{a_i}$ is a factor of $p-1$ and $p_i^{f_i}$ is a factor of $k$, the difference in the exponent is:
    $$\frac{p-1}{p_{i}} = \frac{p_{1}^{a_{1}} \cdots p_{i}^{a_{i}-1} \cdots p_{m}^{a_{m}}}{p_{i}}$$
    
    Since $f_i \le a_i - 1$, we have $k \mid \frac{p-1}{p_{i}}$.
    
    By the property of the order, $k \mid \frac{p-1}{p_{i}}$ implies $g^{(p-1)/p_{i}} \equiv 1 \pmod{p}$.
    This contradicts our initial assumption that $g^{(p-1)/p_{i}} \not\equiv 1 \pmod{p}$ for all $i$.
    
    Therefore, the assumption $k < p-1$ is false, meaning $k = p-1$. Hence, $g$ is a primitive root. \qedhere
\end{proof}

\begin{center}
    \vspace{10pt}
    \line(1,0){450}
    \vspace{25pt}
\end{center}

\begin{algorithm}[Finding a Primitive Root]
    Let $p>2$ be prime and $p-1 = \prod_{i=1}^{m}p_{i}^{a_{i}}$ be the prime factorization of $p-1$, with $p_{i}$ distinct primes and $a_{i} \geq 1$ for $i=1, \dots, m$.
    \begin{enumerate}
        \item Start with $g=2$.
        \item If $g^{(p-1)/p_{i}} \not\equiv 1 \pmod{p}$ for all $i=1, \dots, m$, then $g$ is a primitive root.
        \item If not, increment $g \rightarrow g+1$ and return to Step 2.
    \end{enumerate}
\end{algorithm}